\documentclass[10pt,a4paper]{amsart}
\usepackage{amsmath,amssymb,amsthm,graphicx}
\usepackage{verbatim}
\newenvironment{reference}{\comment}{\endcomment}
%\newenvironment{reference}{}{}
\newenvironment{slogan}{\comment}{\endcomment}
\newenvironment{history}{\comment}{\endcomment}

% For commutative diagrams we use Xy-pic
\usepackage[all]{xy}

% We use 2cell for 2-commutative diagrams.
\xyoption{2cell}
% Theorem environments.
%
\theoremstyle{plain}
\newtheorem{theorem}[subsection]{Theorem}
\newtheorem{proposition}[subsection]{Proposition}
\newtheorem{lemma}[subsection]{Lemma}

\theoremstyle{definition}
\newtheorem{definition}[subsection]{Definition}
\newtheorem{example}[subsection]{Example}
\newtheorem{exercise}[subsection]{Exercise}
\newtheorem{situation}[subsection]{Situation}

\theoremstyle{remark}
\newtheorem{remark}[subsection]{Remark}
\newtheorem{remarks}[subsection]{Remarks}

\numberwithin{equation}{subsection}

% Macros
%
\def\lim{\mathop{\mathrm{lim}}\nolimits}
\def\colim{\mathop{\mathrm{colim}}\nolimits}
\def\Spec{\mathop{\mathrm{Spec}}}
\def\Hom{\mathop{\mathrm{Hom}}\nolimits}
\def\Ext{\mathop{\mathrm{Ext}}\nolimits}
\def\SheafHom{\mathop{\mathcal{H}\!\mathit{om}}\nolimits}
\def\SheafExt{\mathop{\mathcal{E}\!\mathit{xt}}\nolimits}
\def\Sch{\mathit{Sch}}
\def\Mor{\mathop{\mathrm{Mor}}\nolimits}
\def\Ob{\mathop{\mathrm{Ob}}\nolimits}
\def\Sh{\mathop{\mathit{Sh}}\nolimits}
\def\NL{\mathop{N\!L}\nolimits}
\def\CH{\mathop{\mathrm{CH}}\nolimits}
\def\proetale{{pro\text{-}\acute{e}tale}}
\def\etale{{\acute{e}tale}}
\def\QCoh{\mathit{QCoh}}
\def\Ker{\mathop{\mathrm{Ker}}}
\def\Im{\mathop{\mathrm{Im}}}
\def\Coker{\mathop{\mathrm{Coker}}}
\def\Coim{\mathop{\mathrm{Coim}}}

% Boxtimes
%
\DeclareMathSymbol{\boxtimes}{\mathbin}{AMSa}{"02}

%
% Macros for moduli stacks/spaces
%
\def\QCohstack{\mathcal{QC}\!\mathit{oh}}
\def\Cohstack{\mathcal{C}\!\mathit{oh}}
\def\Spacesstack{\mathcal{S}\!\mathit{paces}}
\def\Quotfunctor{\mathrm{Quot}}
\def\Hilbfunctor{\mathrm{Hilb}}
\def\Curvesstack{\mathcal{C}\!\mathit{urves}}
\def\Polarizedstack{\mathcal{P}\!\mathit{olarized}}
\def\Complexesstack{\mathcal{C}\!\mathit{omplexes}}
% \Pic is the operator that assigns to X its picard group, usage \Pic(X)
% \Picardstack_{X/B} denotes the Picard stack of X over B
% \Picardfunctor_{X/B} denotes the Picard functor of X over B
\def\Pic{\mathop{\mathrm{Pic}}\nolimits}
\def\Picardstack{\mathcal{P}\!\mathit{ic}}
\def\Picardfunctor{\mathrm{Pic}}
\def\Deformationcategory{\mathcal{D}\!\mathit{ef}}
\usepackage{cleveref}
\crefname{originallemma}{original lemma}{original lemmas}
\Crefname{originallemma}{Original lemma}{Original lemmas}
\makeatletter
\newlabel{lemma-duality-proper-over-field}{{0FVV}{source}{}{sourcepointer-1}{}}
\newlabel{lemma-duality-proper-over-field@cref}{{[originallemma][0][]0FVV}{[1][0][]source}{}{}{}}
\makeatother
\crefname{originallemma}{original lemma}{original lemmas}
\Crefname{originallemma}{Original lemma}{Original lemmas}
\makeatletter
\newlabel{lemma-twisted-inverse-image}{{0A9E}{source}{}{sourcepointer-2}{}}
\newlabel{lemma-twisted-inverse-image@cref}{{[originallemma][0][]0A9E}{[1][0][]source}{}{}{}}
\makeatother
\crefname{originallemma}{original lemma}{original lemmas}
\Crefname{originallemma}{Original lemma}{Original lemmas}
\makeatletter
\newlabel{coherent-lemma-proper-over-affine-cohomology-finite}{{02O6}{source}{}{sourcepointer-3}{}}
\newlabel{coherent-lemma-proper-over-affine-cohomology-finite@cref}{{[originallemma][0][]02O6}{[1][0][]source}{}{}{}}
\makeatother
\crefname{originalsection}{original section}{original sections}
\Crefname{originalsection}{Original section}{Original sections}
\makeatletter
\newlabel{example-affine-twisted-inverse-image}{{0A9F}{source}{}{sourcepointer-4}{}}
\newlabel{example-affine-twisted-inverse-image@cref}{{[originalsection][0][]0A9F}{[1][0][]source}{}{}{}}
\makeatother
\crefname{originallemma}{original lemma}{original lemmas}
\Crefname{originallemma}{Original lemma}{Original lemmas}
\makeatletter
\newlabel{lemma-more-base-change}{{0AA8}{source}{}{sourcepointer-5}{}}
\newlabel{lemma-more-base-change@cref}{{[originallemma][0][]0AA8}{[1][0][]source}{}{}{}}
\makeatother
\crefname{originalreference}{original reference}{original references}
\Crefname{originalreference}{Original reference}{Original references}
\makeatletter
\newlabel{remark-relative-dualizing-complex}{{0B6S}{source}{}{sourcepointer-6}{}}
\newlabel{remark-relative-dualizing-complex@cref}{{[originalreference][0][]0B6S}{[1][0][]source}{}{}{}}
\makeatother
\title{The dualizing complex of a proper scheme over a field is independent of the field of definition and compatible with field extensions}
\author{The Stacks Project Authors}
\date{stacks-project:duality.tex:local-snapshot-20261012; GNU FDL 1.2 or later}
\begin{document}
\maketitle
\noindent\textit{Editorial note.} All delivered mathematical text is the author's own, copied from the pinned Stacks Project chapter file duality.tex; no statement or proof step was rewritten, reordered or removed.\par\medskip
\noindent\textit{Editorial note.} Preamble provenance: the pinned chapter file opens with an input of the shared preamble and carries no preamble text of its own, so the project's single Stacks-wide preamble staged at sources/stacks-duality/preamble.tex is excerpted verbatim under the recipe fixed by LANE\_BRIEF 22.1, as three spans, lines 8-13 (verbatim package and the reference/slogan/history environments), lines 14-18 (Xy-pic setup) and lines 154-217 (the native theorem declarations and mathematical macros). Nothing was added, removed or redefined; the wrapper class and the amsmath/amssymb/amsthm/graphicx base packages are supplied by the assembler.\par\medskip
\noindent\textit{Editorial note.} The counted statement and proof contain no cross reference, no citation and no included file, so the delivered body is self contained.\par\medskip
\noindent\textit{Editorial note.} Canonical location of the counted proposition: Stacks tag 0FW1 (sources/stacks-duality/pinned-original-source.tex line 7163).\par\medskip
\noindent\textit{Allowed standard prerequisites (editorial).}\begin{itemize}
\item The dualizing complex of a proper scheme over a field and the description of the adjoint on affines. Same chapter, Lemma labelled lemma-duality-proper-over-field and Example labelled example-affine-twisted-inverse-image, cited at proof lines 7207-7219 and 7229-7237. Fixes the objects and the affine computation of the shriek of a field.
\item Finiteness of the field extension arising from the factorisation. Cohomology of Schemes, Lemma labelled coherent-lemma-proper-over-affine-cohomology-finite, cited at proof lines 7221-7223. Makes k'/k finite.
\item Base change for the right adjoint along a flat morphism. Same chapter, Lemma labelled lemma-more-base-change and Remark labelled remark-relative-dualizing-complex, cited at proof lines 7239-7240. Gives part (2).
\end{itemize}\medskip
\ifdefined\hypertarget\hypertarget{sourcepointer-1}{}\fi
\noindent\textit{Source reference.} Original lemma 0FVV: sources/stacks-duality/pinned-original-source.tex, source lines 6934--6934. Same-chapter pointer into the author environment lemma at source line 6934; invoked at proof line 7181 of the counted proof. It is retained only as a navigation pointer; every load-bearing step is inside the delivered span.\par\medskip
\ifdefined\hypertarget\hypertarget{sourcepointer-2}{}\fi
\noindent\textit{Source reference.} Original lemma 0A9E: sources/stacks-duality/pinned-original-source.tex, source lines 653--653. Same-chapter pointer into the author environment lemma at source line 653; invoked at proof line 7184 of the counted proof. It is retained only as a navigation pointer; every load-bearing step is inside the delivered span.\par\medskip
\ifdefined\hypertarget\hypertarget{sourcepointer-3}{}\fi
\noindent\textit{Source reference.} Original lemma 02O6: sources/stacks-chow/tags.txt, source lines 3505--3505. Cross-chapter standard prerequisite located in the project tag list staged in the local snapshot; invoked at proof line 7191 of the counted proof. It is retained only as a navigation pointer; every load-bearing step is inside the delivered span.\par\medskip
\ifdefined\hypertarget\hypertarget{sourcepointer-4}{}\fi
\noindent\textit{Source reference.} Original section 0A9F: sources/stacks-duality/pinned-original-source.tex, source lines 681--681. Same-chapter pointer into the author environment example at source line 681; invoked at proof line 7197 of the counted proof. It is retained only as a navigation pointer; every load-bearing step is inside the delivered span.\par\medskip
\ifdefined\hypertarget\hypertarget{sourcepointer-5}{}\fi
\noindent\textit{Source reference.} Original lemma 0AA8: sources/stacks-duality/pinned-original-source.tex, source lines 1636--1636. Same-chapter pointer into the author environment lemma at source line 1636; invoked at proof line 7203 of the counted proof. It is retained only as a navigation pointer; every load-bearing step is inside the delivered span.\par\medskip
\ifdefined\hypertarget\hypertarget{sourcepointer-6}{}\fi
\noindent\textit{Source reference.} Original reference 0B6S: sources/stacks-duality/pinned-original-source.tex, source lines 2945--2945. Same-chapter pointer into the author environment remark at source line 2945; invoked at proof line 7204 of the counted proof. It is retained only as a navigation pointer; every load-bearing step is inside the delivered span.\par\medskip
\begin{lemma}
\label{lemma-sanity-check-duality}
Let $X$ be a proper scheme over a field $k$. Let $\omega_X^\bullet$ and
$\omega_X$ be as in Lemma \ref{lemma-duality-proper-over-field}.
\begin{enumerate}
\item If $X \to \Spec(k)$ factors as $X \to \Spec(k') \to \Spec(k)$
for some field $k'$, then $\omega_X^\bullet$ and $\omega_X$
are as in Lemma \ref{lemma-duality-proper-over-field} for the morphism
$X \to \Spec(k')$.
\item If $K/k$ is a field extension, then the pullback of
$\omega_X^\bullet$ and $\omega_X$ to the base change $X_K$
are as in  Lemma \ref{lemma-duality-proper-over-field} for the morphism
$X_K \to \Spec(K)$.
\end{enumerate}
\end{lemma}

\begin{proof}
Denote $f : X \to \Spec(k)$ the structure morphism and denote
$f' : X \to \Spec(k')$ the given factorization.
In the proof of Lemma \ref{lemma-duality-proper-over-field}
we took $\omega_X^\bullet = a(\mathcal{O}_{\Spec(k)})$
where $a$ be is the right adjoint of
Lemma \ref{lemma-twisted-inverse-image} for $f$.
Thus we have to show
$a(\mathcal{O}_{\Spec(k)}) \cong a'(\mathcal{O}_{\Spec(k)})$
where $a'$ be is the right adjoint of
Lemma \ref{lemma-twisted-inverse-image} for $f'$.
Since $k' \subset H^0(X, \mathcal{O}_X)$ we see that $k'/k$ is a finite
extension (Cohomology of Schemes, Lemma
\ref{coherent-lemma-proper-over-affine-cohomology-finite}).
By uniqueness of adjoints we have $a = a' \circ b$ where
$b$ is the right adjoint of Lemma
\ref{lemma-twisted-inverse-image} for $g : \Spec(k') \to \Spec(k)$.
Another way to say this: we have $f^! = (f')^! \circ g^!$.
Thus it suffices to show that $\Hom_k(k', k) \cong k'$ as
$k'$-modules, see Example \ref{example-affine-twisted-inverse-image}.
This holds because these are $k'$-vector spaces of
the same dimension (namely dimension $1$).

\medskip\noindent
Proof of (2). This holds because we have base change for $a$ by
Lemma \ref{lemma-more-base-change}. See discussion in
Remark \ref{remark-relative-dualizing-complex}.
\end{proof}

\end{document}
